% Capítulo: La preexistencia de Dios (generado desde 01_Preexistencia/01_Preexistencia_capitulo.tex)
% Texto ancla: Éxodo 3:13–15. Clasificación: Ser; incomunicable.

\section{Antes que todo}

Toda criatura tiene un antes. Los montes tuvieron un día en que no existían; el sol, el mar y el ser humano también. La pregunta de este capítulo es la más antigua que puede hacerse: ¿qué había antes de todo? La Escritura responde con un Nombre. Antes de que hubiera montes, tierra o mundo, ya estaba Dios:

\hebverso{בְּטֶרֶם הָרִים יֻלָּדוּ וַתְּחוֹלֵל אֶרֶץ וְתֵבֵל \\ וּמֵעוֹלָם עַד־עוֹלָם אַתָּה אֵל}
\begin{verso}{Salmo 90:2}
Antes que naciesen los montes Y formases la tierra y el mundo, Desde el siglo y hasta el siglo, tú eres Dios.
\end{verso}

El material de la serie define el atributo en una línea: \destacar{Dios existe por sí mismo y no depende de nada}. Esa frase contiene tres afirmaciones, y el capítulo las sigue en ese orden:

\begin{center}\small\renewcommand{\arraystretch}{1.25}
\begin{tabular}{@{}L{3.2cm}L{6.4cm}L{4.8cm}@{}}
\toprule
\textbf{Rasgo} & \textbf{Pregunta} & \textbf{Textos centrales} \\ \midrule
Preexistencia & ¿Qué había antes de todo? Dios. & Sal 90:2; Is 44:6; Jn 1:1 \\
Existencia por sí mismo & ¿De dónde viene Dios? De nadie: es. & Éx 3:14; Jn 5:26 \\
Autosuficiencia & ¿Qué le falta a Dios? Nada. & Sal 50:12; Job 41:11; Hch 17:25 \\
\bottomrule
\end{tabular}\par
{\footnotesize\color{gris}\textit{Tabla 1.} Los tres rasgos del atributo. \textit{Nota.} Elaboración propia a partir de la definición del material de la serie.}
\end{center}

El verso que une esta obra pide gloriarse en ``entenderme y conocerme, que yo soy \SENOR'' (Jer 9:24 [9:23]). El primer paso para conocer a alguien es saber de dónde viene. Con Dios esa pregunta se responde de un modo único: no viene de ningún lugar ni de nadie. Todo lo demás viene de él.

%===================================================================
\section{Texto ancla: Éxodo 3:13--15}

\subsection{Una zarza que no se consume}
Moisés apacienta las ovejas de su suegro en Horeb, ``monte de Dios'', cuando ve algo que no puede explicar:

\hebverso{וְהִנֵּה הַסְּנֶה בֹּעֵר בָּאֵשׁ וְהַסְּנֶה אֵינֶנּוּ אֻכָּל}
\begin{verso}{Éxodo 3:2}
Y se le apareció el Angel de \SENOR{} en una llama de fuego en medio de una zarza; y él miró, y vio que la zarza ardía en fuego, y la zarza no se consumía.
\end{verso}

Todo fuego de este mundo depende de algo: necesita leña y se apaga cuando la consume. Este fuego arde y la zarza sigue entera. Moisés se acerca a ver ``por qué causa la zarza no se quema'' (3:3), y desde ese fuego Dios le habla.

\precision{leer la zarza como una imagen de un Dios que no depende de combustible es una lectura propia de esta serie. El texto subraya el asombro de Moisés y la santidad del lugar (3:5); la tradición judía leyó la zarza de otras maneras, por ejemplo como figura de Israel que sufre en Egipto sin ser destruido.}

\subsection{``¿Cuál es su nombre?''}
Dios envía a Moisés a Faraón y le da una primera promesa: ``yo estaré contigo'' (3:12). Moisés anticipa la pregunta del pueblo:

\begin{verso}{Éxodo 3:13}
Dijo Moisés a Dios: He aquí que llego yo a los hijos de Israel, y les digo: El Dios de vuestros padres me ha enviado a vosotros. Si ellos me preguntaren: ¿Cuál es su nombre?, ¿qué les responderé?
\end{verso}

En Israel el nombre dice quién es alguien. Un esclavo en Egipto conocía los nombres de muchos dioses, cada uno con su templo, su ciudad y su función. Moisés pide el Nombre que distinga a este Dios de todos ellos.

\subsection{``YO SOY EL QUE SOY''}
\hebverso{וַיֹּאמֶר אֱלֹהִים אֶל־מֹשֶׁה אֶהְיֶה אֲשֶׁר אֶהְיֶה \\ וַיֹּאמֶר כֹּה תֹאמַר לִבְנֵי יִשְׂרָאֵל אֶהְיֶה שְׁלָחַנִי אֲלֵיכֶם}
\begin{verso}{Éxodo 3:14}
Y respondió Dios a Moisés: YO SOY EL QUE SOY. Y dijo: Así dirás a los hijos de Israel: YO SOY me envió a vosotros.
\end{verso}

\begin{filologia}{Éx 3:14}
\heb{אֶהְיֶה} \emph{ehyeh} & \heb{היה} & primera persona del imperfecto (qal) de \emph{hayá}, ``ser, existir, llegar a ser, estar presente''. El imperfecto hebreo no marca un tiempo fijo: puede traducirse ``soy'', ``seré'' o ``estaré''. \\
\heb{אֲשֶׁר} \emph{asher} & --- & relativo: ``el que'', ``lo que''. La fórmula se cierra sobre sí misma: Dios se define por Dios. \\
\heb{אֶהְיֶה שְׁלָחַנִי} \emph{ehyeh shelajani} & \heb{שׁלח} & ``\emph{Ehyeh} me envió'': el verbo se vuelve Nombre propio. \\
\grc{ἐγώ εἰμι ὁ ὤν} & \grc{εἰμί} & la Septuaginta traduce ``Yo soy el que es'' (\grc{ὁ ὤν}, participio de ``ser''), y luego ``\grc{ὁ ὤν} me envió''. \\
\end{filologia}

La misma forma \emph{ehyeh} ya había aparecido dos versículos antes: \heb{כִּי־אֶהְיֶה עִמָּךְ}, ``yo estaré contigo'' (3:12). El Dios que \emph{es} por sí mismo es también el Dios que \emph{está} con su siervo. Las dos lecturas del verbo (existencia y presencia) no se excluyen en el texto: el Dios que no depende de nadie es el que baja a librar a un pueblo que depende de todo (3:7--8).

\precision{la traducción de \emph{ehyeh asher ehyeh} se discute. ``Yo soy el que soy'' subraya el ser (así la LXX y la RVR1960); ``Seré el que seré'' o ``Estaré como el que estaré'' subraya la presencia fiel en la historia, apoyada en 3:12 (HALOT). Esta serie sostiene que el hebreo permite las dos lecturas y que el contexto las une. La lectura judía tradicional, como se verá en la sección~\ref{sec:pre-judaismo}, insiste en la presencia.}

\subsection{El Nombre para siempre}
\hebverso{זֶה־שְּׁמִי לְעֹלָם וְזֶה זִכְרִי לְדֹר דֹּר}
\begin{verso}{Éxodo 3:15}
Además dijo Dios a Moisés: Así dirás a los hijos de Israel: \SENOR, el Dios de vuestros padres, el Dios de Abraham, Dios de Isaac y Dios de Jacob, me ha enviado a vosotros. Este es mi nombre para siempre; con él se me recordará por todos los siglos.
\end{verso}

En el v. 14 Dios dice \emph{Ehyeh}, ``Yo soy''; en el v. 15 da el Nombre \heb{יהוה}, que Israel pronunciará en tercera persona. El texto pone uno junto al otro: el tetragrama se escucha como ``Él es''. Y lo vincula a Abraham, Isaac y Jacob: el Dios que existía antes de los patriarcas es el que se ató a ellos con un pacto.

\precision{relacionar \heb{יהוה} con el verbo \emph{hayá} (en su forma antigua \emph{hawá}) es la explicación que sugiere el propio texto al poner Éx 3:14 y 3:15 juntos; el origen exacto del Nombre se discute en la investigación (HALOT). La pronunciación original no se conserva; la tradición judía lee \emph{Adonay}.}

%===================================================================
\section{La palabra: decir ``antes'' de Dios}

\subsection{Una idea sin sustantivo}
La palabra ``preexistencia'' no aparece en la Biblia. El hebreo no forma sustantivos abstractos para Dios: lo dice con verbos, pronombres y adverbios de tiempo. El atributo se reconoce por un conjunto de expresiones que se repiten:

\begin{center}\small\renewcommand{\arraystretch}{1.16}
\begin{tabular}{@{}L{3.7cm}L{1.0cm}L{5.4cm}L{4.3cm}@{}}
\toprule
\textbf{Expresión} & \textbf{Veces} & \textbf{Sentido y traducción RVR1960} & \textbf{Textos} \\ \midrule
\heb{אֶהְיֶה} \emph{ehyeh} & 51 & ``seré, estaré''; en Éx 3:14, ``YO SOY'' & Éx 3:12, 14; Pr 8:30 \\
\heb{אֲנִי הוּא} \emph{ani hu} & 8 & ``yo soy'', ``yo mismo'', ``yo era'' & Dt 32:39; Is 41:4; 43:10, 13; 46:4; 48:12 \\
\heb{בְּטֶרֶם} \emph{béterem} & 45 & ``antes que'' (pocas veces de Dios) & Sal 90:2; Pr 8:25; Jer 1:5 \\
\heb{מִקֶּדֶם} \emph{miqqédem}, \heb{אֱלֹהֵי קֶדֶם} & 24; 1 & ``desde el principio'', ``desde la antigüedad''; ``el eterno Dios'' & Dt 33:27; Is 46:10; Mi 5:2; Hab 1:12 \\
\heb{רִאשׁוֹן} y \heb{אַחֲרוֹן} & 3 & ``el primero\ldots{} el postrero'' & Is 41:4; 44:6; 48:12 \\
\heb{מֵעוֹלָם עַד־עוֹלָם} & 6 & ``desde el siglo y hasta el siglo'', ``desde la eternidad hasta la eternidad'' & Sal 90:2; 103:17; Neh 9:5 \\ \midrule
\grc{πρὸ καταβολῆς κόσμου} & 3 & ``antes de la fundación del mundo'' & Jn 17:24; Ef 1:4; 1~P 1:20 \\
\grc{ἀπὸ καταβολῆς κόσμου} & 6 & ``desde la fundación del mundo'' & Mt 25:34; Ap 13:8 \\
\grc{πρὸ χρόνων αἰωνίων}; \grc{πρὸ τῶν αἰώνων} & 2; 1 & ``antes de los tiempos de los siglos''; ``antes de los siglos'' & 2~Ti 1:9; Tit 1:2; 1~Co 2:7 \\
\grc{πρὸ πάντων} & 1 & ``antes de todas las cosas'' (de Cristo) & Col 1:17 \\
\grc{πρὶν Ἀβραὰμ γενέσθαι} & 1 & ``antes que Abraham fuese'' & Jn 8:58 \\
\grc{ἐγώ εἰμι} absoluto & 7 & ``yo soy'', sin predicado & Jn 8:24, 28, 58; 13:19; 18:5, 6, 8 \\
\grc{προγινώσκω}; \grc{πρόγνωσις} & 5; 2 & ``conocer antes''; ``anticipado conocimiento'' & Hch 2:23; Ro 8:29; 1~P 1:20 \\
\grc{προσδέομαι} & 1 & ``necesitar de algo'' (negado de Dios) & Hch 17:25 \\
\bottomrule
\end{tabular}\par
{\footnotesize\color{gris}\textit{Tabla 2.} El espectro de la preexistencia. \textit{Nota.} Conteo propio con la BHS y NA27 con números Strong; la RVR1960 sin títulos de sección. En \emph{miqqédem} el conteo incluye el sentido local ``al oriente''; en \emph{ani hu}, 1~Cr 21:17, donde lo dice David.}
\end{center}

\precision{los conteos son propios y conviene verificarlos en una concordancia (e-Sword). Las cifras de \emph{ehyeh} y \emph{béterem} cuentan todas las apariciones, aunque la mayoría no habla de Dios; se dan para mostrar que el lenguaje es ordinario y que lo extraordinario es el sujeto. En \emph{ego eimí} solo se cuentan los usos absolutos de Juan; con predicado (``yo soy el pan'', ``la luz'', ``la puerta'') la fórmula aparece muchas más veces.}

\subsection{\emph{Ani hu}: ``Yo soy''}
La expresión más cercana a Éx 3:14 es \heb{אֲנִי הוּא}, literalmente ``yo (soy) él''. Dios la usa en el Cántico de Moisés y en Isaías 40--48, siempre frente a los dioses que no existían antes ni existirán después:

\hebverso{כִּי־אֲנִי הוּא לְפָנַי לֹא־נוֹצַר אֵל וְאַחֲרַי לֹא יִהְיֶה}
\begin{verso}{Isaías 43:10}
Vosotros sois mis testigos, dice \SENOR, y mi siervo que yo escogí, para que me conozcáis y creáis, y entendáis que yo mismo soy; antes de mí no fue formado dios, ni lo será después de mí.
\end{verso}

\begin{filologia}{Is 43:10 y 43:13}
\heb{אֲנִי הוּא} \emph{ani hu} & --- & ``yo (soy) él'': frase nominal sin verbo. La LXX la traduce \grc{ἐγώ εἰμι}. \\
\heb{לְפָנַי לֹא־נוֹצַר אֵל} & \heb{יצר} & ``antes de mí no fue formado dios'': \emph{yatsar} es el verbo del alfarero (Gn 2:7). Todo dios ``formado'' tiene un antes; Dios no. \\
\heb{גַּם־מִיּוֹם אֲנִי הוּא} & \heb{יום} & ``aun antes que hubiera día, yo era'' (43:13, RVR1960); literalmente ``también desde el día, yo soy él''. \\
\end{filologia}

Israel es llamado ``testigo'' de este Dios: lo conoce y cree en él (\heb{וְתַאֲמִינוּ}, \emph{veta'aminu}, de la raíz \emph{amán}). En la Septuaginta, la frase \grc{ἐγώ εἰμι} de estos textos será la que Jesús pronuncie en el Evangelio de Juan (sección~\ref{sec:pre-abraham}).

\begin{center}
\begin{tikzpicture}[font=\small,node distance=0.9cm]
 \node[draw=azul,line width=0.8pt,fill=papiro,rounded corners=3pt,text width=3.6cm,align=center,inner sep=6pt] (h) {\textbf{Hebreo}\\ \heb{אֶהְיֶה} (Éx 3:14)\\ \heb{אֲנִי הוּא} (Dt 32:39; Is 43:10)};
 \node[draw=azul,line width=0.8pt,fill=papiro,rounded corners=3pt,text width=3.6cm,align=center,inner sep=6pt,right=of h] (l) {\textbf{Septuaginta}\\ \grc{ἐγώ εἰμι ὁ ὤν}\\ \grc{ἐγώ εἰμι}};
 \node[draw=dorado,line width=1.6pt,fill=azul,text=papiro,rounded corners=3pt,text width=3.6cm,align=center,inner sep=6pt,right=of l] (p) {\textbf{Jesús}\\ \grc{πρὶν Ἀβραὰμ γενέσθαι}\\ \grc{ἐγὼ εἰμί} (Jn 8:58)};
 \draw[-{Latex[length=3mm]},dorado,line width=1.4pt] (h)--(l);
 \draw[-{Latex[length=3mm]},dorado,line width=1.4pt] (l)--(p);
\end{tikzpicture}\par
{\footnotesize\color{gris}\textit{Figura 1.} Del Nombre a Jesús. \textit{Nota.} Elaboración propia.}
\end{center}

\subsection{\emph{Béterem}: antes que}
\heb{בְּטֶרֶם} (\emph{béterem}) es una palabra común: ``antes de que'' llegara alguien, ``antes de que'' se acostaran (Gn 19:4). Aplicada a Dios, abre un tiempo en que solo él estaba. Sal 90:2 la usa para Dios mismo: antes de los montes, ``tú eres Dios''. El salmista no escribe ``tú eras'': \heb{אַתָּה אֵל} (\emph{atá El}) no tiene verbo, como si Dios no tuviera pasado. Jeremías la usa para el plan de Dios sobre una persona:

\hebverso{בְּטֶרֶם אֶצָּרְךָ בַבֶּטֶן יְדַעְתִּיךָ}
\begin{verso}{Jeremías 1:5}
Antes que te formase en el vientre te conocí, y antes que nacieses te santifiqué, te di por profeta a las naciones.
\end{verso}

\begin{filologia}{Jer 1:5}
\heb{בְּטֶרֶם} \emph{béterem} & --- & ``antes que'': el conocimiento de Dios precede a la existencia del profeta. \\
\heb{אֶצָּרְךָ} \emph{etsorjá} & \heb{יצר} & ``te formase'': el verbo del alfarero (Gn 2:7; Is 43:10). \\
\heb{יְדַעְתִּיךָ} \emph{yedatija} & \heb{ידע} & ``te conocí'': conocer en hebreo es relación y elección, no solo información (Gn 18:19; Am 3:2). \\
\end{filologia}

El mismo libro que dice ``conocerme'' (Jer 9:24) comienza diciendo que Dios conoció primero. Aquí la preexistencia de Dios toca la \destacar{presciencia}: Dios conoce \emph{antes} porque \emph{es} antes.

\subsection{\emph{Miqqédem}: desde siempre}
\heb{קֶדֶם} (\emph{qédem}) significa ``lo que está delante'': el oriente, por donde sale el sol, y también lo antiguo, lo que está ``antes'' en el tiempo. Moisés llama a Dios \heb{אֱלֹהֵי קֶדֶם}, ``el eterno Dios'' (Dt 33:27), y Habacuc pregunta en medio de la crisis:

\hebverso{הֲלוֹא אַתָּה מִקֶּדֶם יְהוָה אֱלֹהַי קְדֹשִׁי לֹא נָמוּת}
\begin{verso}{Habacuc 1:12}
¿No eres tú desde el principio, oh \SENOR, Dios mío, Santo mío? No moriremos. Oh \SENOR, para juicio lo pusiste; y tú, oh Roca, lo fundaste para castigar.
\end{verso}

Habacuc saca una consecuencia: si Dios es ``desde el principio'', su pueblo ``no morirá''. La preexistencia de Dios no es un dato frío: es el fundamento de la esperanza de quien depende de él.

\subsection{``¿Quién me ha dado a mí primero?''}
El rasgo de la autosuficiencia tiene también su vocabulario. Dios pregunta en Job:

\hebverso{מִי הִקְדִּימַנִי וַאֲשַׁלֵּם תַּחַת כָּל־הַשָּׁמַיִם לִי־הוּא}
\begin{verso}{Job 41:11 [41:3]}
¿Quién me ha dado a mí primero, para que yo restituya? Todo lo que hay debajo del cielo es mío.
\end{verso}

\hebverso{אִם־אֶרְעַב לֹא־אֹמַר לָךְ כִּי־לִי תֵבֵל וּמְלֹאָהּ}
\begin{verso}{Salmo 50:12}
Si yo tuviese hambre, no te lo diría a ti; Porque mío es el mundo y su plenitud.
\end{verso}

\begin{filologia}{Job 41:11 y Sal 50:12}
\heb{הִקְדִּימַנִי} \emph{hiqdimani} & \heb{קדם} & ``se me adelantó'', ``me dio primero'': hifil de la misma raíz de \emph{qédem}. Nadie puede llegar \emph{antes} que Dios para dejarlo en deuda. \\
\heb{וַאֲשַׁלֵּם} \emph{va'ashalém} & \heb{שׁלם} & ``para que yo pague, restituya''. \\
\heb{אִם־אֶרְעַב} \emph{im er'av} & \heb{רעב} & ``si tuviera hambre'': condición irreal. Los sacrificios no alimentan a Dios (50:13). \\
\end{filologia}

Job 41:11 une los dos rasgos con una sola raíz: Dios es \emph{qédem} (antes) y por eso nadie se le \emph{adelanta} (\emph{hiqdim}). Pablo citará este versículo en Ro 11:35 (sección~\ref{sec:pre-pablo}).

%===================================================================
\section{El Nombre que revela la preexistencia}

En la Escritura el nombre de Dios revela lo que él es. Moisés pidió el Nombre en la zarza y recibió \emph{Ehyeh}, ``YO SOY'' (Éx 3:14), el texto ancla de este capítulo, que el libro estudiará también desde la eternidad de Dios. Pero hay un título que Dios usa para decir, en particular, que está antes de todo: \destacar{``el primero''}.

\subsection{``Yo soy el primero, y yo soy el postrero''}
\hebverso{כֹּה־אָמַר יְהוָה מֶלֶךְ־יִשְׂרָאֵל וְגֹאֲלוֹ יְהוָה צְבָאוֹת \\ אֲנִי רִאשׁוֹן וַאֲנִי אַחֲרוֹן וּמִבַּלְעָדַי אֵין אֱלֹהִים}
\begin{verso}{Isaías 44:6}
Así dice \SENOR{} Rey de Israel, y su Redentor, \SENOR{} de los ejércitos: Yo soy el primero, y yo soy el postrero, y fuera de mí no hay Dios.
\end{verso}

\begin{filologia}{Is 44:6}
\heb{אֲנִי רִאשׁוֹן} \emph{ani rishón} & \heb{ראשׁ} & ``yo (soy) primero''. \emph{Rishón} viene de \emph{rosh}, ``cabeza'': el que encabeza. De la misma raíz es \emph{reshit}, ``principio'' (Gn 1:1). \\
\heb{וַאֲנִי אַחֲרוֹן} \emph{va'ani ajarón} & \heb{אחר} & ``y yo (soy) último'': el que queda cuando todo lo demás pasó. \\
\heb{וּמִבַּלְעָדַי אֵין אֱלֹהִים} & \heb{בלעד} & ``y fuera de mí no hay Dios'': el título es exclusivo. \\
\heb{גֹּאֲלוֹ} \emph{go'aló} & \heb{גאל} & ``su Redentor'': el Primero es también el pariente que rescata. \\
\end{filologia}

Isaías repite el título tres veces, siempre frente a los ídolos y al exilio: ``Yo \SENOR, el primero, y yo mismo con los postreros'' (41:4); ``Yo mismo, yo el primero, yo también el postrero'' (48:12). Dios no es el primero de una lista de seres que vinieron después de él al mismo nivel; es el primero porque todo lo demás viene de él, y el último porque nada lo sobrevive. Y lo une con ``Redentor'': el que está antes de todo se ató a Israel para rescatarlo.

\subsection{``El Alfa y la Omega''}
El Apocalipsis traduce el título al alfabeto griego. Dios lo dice de sí mismo:

\begin{verso}{Apocalipsis 1:8}
Yo soy el Alfa y la Omega, principio y fin, dice el Señor, el que es y que era y que ha de venir, el Todopoderoso.
\end{verso}

\noindent y Cristo lo dice de sí mismo con las palabras de Isaías:

\begin{verso}{Apocalipsis 1:17}
Cuando le vi, caí como muerto a sus pies. Y él puso su diestra sobre mí, diciéndome: No temas; yo soy el primero y el último.
\end{verso}
\begin{verso}{Apocalipsis 22:13}
Yo soy el Alfa y la Omega, el principio y el fin, el primero y el último.
\end{verso}

\begin{filologia}{Ap 1:8, 17; 22:13}
\grc{τὸ ἄλφα καὶ τὸ ὦ} & --- & ``el Alfa y la Omega'': la primera y la última letra; todo lo que se puede decir está entre las dos. \\
\grc{ὁ πρῶτος καὶ ὁ ἔσχατος} & \grc{πρῶτος} & ``el primero y el último'': traducción de \heb{רִאשׁוֹן} y \heb{אַחֲרוֹן} de Is 44:6. \\
\grc{ἡ ἀρχὴ καὶ τὸ τέλος} & \grc{ἀρχή} & ``el principio y el fin'': origen y meta. \\
\end{filologia}

Lo que en Isaías es título exclusivo del \SENOR{} (``fuera de mí no hay Dios''), en el Apocalipsis lo dice Jesús. El título del Dios que está antes de todo pasa sin cambios al Hijo.

\precision{el paso del título de Dios (Ap 1:8; 21:6) a Cristo (Ap 1:17; 22:13) lo analiza Bauckham (1993). La tradición rabínica leyó la misma idea en las letras de \emph{emet}: álef, la primera; tav, la última (véase el capítulo sobre la verdad).}

\subsection{Pablo: ``él que es el principio''}
Pablo no usa la fórmula del Apocalipsis, pero da a Cristo el mismo lugar con otra palabra:

\begin{verso}{Colosenses 1:18}
y él es la cabeza del cuerpo que es la iglesia, él que es el principio, el primogénito de entre los muertos, para que en todo tenga la preeminencia.
\end{verso}

\grc{ὅς ἐστιν ἀρχή}: ``que es (el) principio''. Antes ha dicho ``él es antes de todas las cosas'' (1:17). Para Pablo, Cristo es principio de la creación y principio de la nueva creación: el primero antes de todo y el primero en salir de la muerte. Y toda la confesión termina en el Nombre: Dios ``le dio un nombre que es sobre todo nombre'', para que toda lengua confiese ``que Jesucristo es el Señor'' (Flp 2:9--11), con las palabras que Isaías reservaba al \SENOR{} (Is 45:23).

%===================================================================
\section{Definición}

\begin{tcolorbox}[colback=papiro,colframe=azul,boxrule=1pt,arc=3pt,fonttitle=\bfseries,title=Definición]
La preexistencia de Dios es la perfección eterna por la cual Dios existe antes de todas las cosas, por sí mismo y sin depender de nada. Nadie le dio el ser, nadie lo sostiene y nada le falta; todo lo que existe recibe de él vida, aliento y todas las cosas (Éx 3:14; Sal 90:2; Hch 17:24--25).
\end{tcolorbox}

\begin{center}\small\renewcommand{\arraystretch}{1.2}
\begin{tabular}{@{}L{3.4cm}L{6.6cm}L{4.4cm}@{}}
\toprule
\textbf{Rasgo} & \textbf{Qué afirma} & \textbf{Textos} \\ \midrule
Antes de todo & cuando nada existía, Dios era & Sal 90:2; Is 43:10; Jn 1:1 \\
Por sí mismo & su ser no viene de otro: ``YO SOY'' & Éx 3:14; Jn 5:26 \\
Sin necesidad & no necesita alimento, templo ni ayuda & Sal 50:12; Hch 17:25 \\
Sin deuda & nadie le dio primero & Job 41:11; Ro 11:35 \\
Fuente de todo & todo recibe de él el ser & Neh 9:6; Ro 11:36; Ap 4:11 \\
Con un plan previo & conoce y decide antes de crear & Is 46:10; Jer 1:5; Ef 1:4 \\
\bottomrule
\end{tabular}\par
{\footnotesize\color{gris}\textit{Tabla 3.} Rasgos de la preexistencia. \textit{Nota.} Elaboración propia.}
\end{center}

\precision{la teología clásica llama a este atributo \emph{aseidad} (del latín \emph{a se}, ``por sí mismo'') o \emph{independencia} de Dios (Berkhof, 1938; Grudem, 2007). El material de la serie usa ``preexistencia'' y ``autosuficiencia''. En este capítulo los tres términos se usan juntos: preexistencia mira al \emph{antes}, aseidad al \emph{origen} y autosuficiencia a la \emph{necesidad}.}

\textbf{Clasificación.} La preexistencia pertenece a los atributos del \emph{Ser}: responde a ``¿qué es Dios?'', y abre la serie porque los demás atributos la suponen. Un Dios que dependiera de algo tendría santidad, poder o amor prestados. Es un atributo \emph{incomunicable}: ninguna criatura existe por sí misma. Solo el Hijo tiene ``vida en sí mismo'', y la tiene porque el Padre se la dio desde siempre (Jn 5:26).

\subsection{Preexistencia y eternidad}
La preexistencia y la eternidad se tocan, pero no son lo mismo. La eternidad dice que Dios no tiene principio ni fin y que el tiempo no lo mide (Sal 90:4; 102:27); será el tema del capítulo sobre el Dios eterno. La preexistencia mira solo hacia atrás desde la creación: antes del primer día, Dios ya era. Sal 90:2 contiene las dos ideas en un verso: ``antes que naciesen los montes'' (preexistencia), ``desde el siglo y hasta el siglo'' (eternidad).

\subsection{Preexistencia y presciencia}
Si Dios existía antes de todo, su conocimiento y su decisión también: ``anuncio lo por venir desde el principio, y desde la antigüedad lo que aún no era hecho; que digo: Mi consejo permanecerá'' (Is 46:10). La creación no tomó a Dios por sorpresa ni lo obligó a improvisar. Antes de que existiera el mundo, Dios ya conocía a Jeremías (Jer 1:5), ya había escogido a su pueblo ``en él'' (Ef 1:4) y ya había destinado al Cordero (1~P 1:20).

\begin{center}
\begin{tikzpicture}[font=\small]
 \draw[line width=1.6pt,azul,-{Latex[length=3mm]}] (0,0) -- (14.2,0);
 \fill[dorado] (3.6,-0.25) rectangle (3.75,0.25);
 \node[font=\scriptsize\bfseries,text=azul,above] at (3.68,0.3) {creación (Gn 1:1)};
 \draw[decorate,decoration={brace,amplitude=6pt,mirror},line width=0.8pt,azuloscuro] (0.1,-0.4) -- (3.5,-0.4);
 \node[text width=3.4cm,align=center,font=\scriptsize,below] at (1.8,-0.65) {\textbf{Antes de todo}\\ Dios es (Sal 90:2)\\ conoce y elige (Ef 1:4; Jer 1:5)\\ el Padre ama al Hijo (Jn 17:24)};
 \draw[decorate,decoration={brace,amplitude=6pt,mirror},line width=0.8pt,azuloscuro] (3.9,-0.4) -- (14,-0.4);
 \node[text width=9cm,align=center,font=\scriptsize,below] at (8.95,-0.65) {\textbf{Historia}\\ todo existe ``por su voluntad'' (Ap 4:11) y ``en él subsiste'' (Col 1:17)\\ Dios lleva a cabo lo que decidió ``desde el principio'' (Is 46:10)};
\end{tikzpicture}\par
{\footnotesize\color{gris}\textit{Figura 2.} Antes de la creación y en la historia. \textit{Nota.} Elaboración propia.}
\end{center}

%===================================================================
\section{Teología: el Dios que no necesita nada}

La Escritura no empieza explicando a Dios, sino con Dios: ``En el principio creó Dios'' (Gn 1:1). Antes de ese principio no había cielo, ni tierra, ni ángeles, ni universo; solo Dios. La pregunta teológica que abre el atributo es qué clase de existencia era esa: ¿estaba Dios incompleto, a la espera de algo que la creación vendría a darle? La Escritura responde que no. El que se nombra ``YO SOY EL QUE SOY'' (Éx 3:14) no recibe su ser de otro ni lo completa con otro. Como observa Pink, durante la eternidad pasada Dios estaba ``completo, autosuficiente, satisfecho de Sí mismo y sin necesidad de nada'', y lo estaba ``en la unidad de su naturaleza (aunque subsistiendo igualmente en tres personas divinas)'' (Pink, 2020, cap. 1). Esa plenitud no fue soledad: es la gloria y el amor que el Hijo tenía con el Padre ``antes que el mundo fuese'' (Jn 17:5), que este capítulo estudia más adelante. De esa plenitud se siguen tres consecuencias que la Escritura despliega.

\subsection{La creación fue libre}

Si el universo le hubiera sido necesario a Dios, habría existido siempre junto a él, como una parte de su ser. La Escritura, en cambio, presenta la creación como fruto de una decisión: Dios ``hace todas las cosas según el designio de su voluntad'' (Ef 1:11), y los ancianos del Apocalipsis lo confiesan delante del trono: ``por tu voluntad existen y fueron creadas'' (Ap 4:11). Crear no le añadió nada a su gloria esencial, porque él no cambia: ``Porque yo \SENOR{} no cambio'' (Mal 3:6). El mundo existe porque Dios quiso que existiera, no porque Dios lo necesitara.

\subsection{Nada le añade}

Si nada le falta, nada le añade, y esto vale incluso para lo más alto que la criatura puede ofrecer: la adoración. Los levitas del tiempo de Nehemías bendicen el Nombre de Dios y, en la misma frase, confiesan que está ``alto sobre toda bendición y alabanza'' (Neh 9:5). Eliú lleva el argumento al terreno moral: ``Si fueres justo, ¿qué le darás a él?'' (Job 35:7). El salmo pone la misma verdad en boca de Dios: ``Si yo tuviese hambre, no te lo diría a ti; porque mío es el mundo y su plenitud'' (Sal 50:12). Y Pablo la predica en Atenas: Dios no ``es honrado por manos de hombres, como si necesitase de algo'' (Hch 17:25). La adoración no enriquece a Dios; es la respuesta de la criatura a una plenitud que ya existe.

\subsection{Nadie lo pone en deuda}

Si todo lo que la criatura tiene viene de Dios, nadie puede ponerlo bajo obligación. Dios mismo lo dice en Job: ``¿Quién me ha dado a mí primero, para que yo restituya? Todo lo que hay debajo del cielo es mío'' (Job 41:11). Pablo cierra con ese versículo su exposición del plan de Dios con Israel y con las naciones: ``¿O quién le dio a él primero, para que le fuese recompensado?'' (Ro 11:35). Lo que el hombre da a Dios ya era de Dios; por eso la relación con él no puede fundarse en el mérito, sino en la gracia. Isaías lleva la comparación a la escala de la historia:

\begin{verso}{Isaías 40:15–18}
He aquí que las naciones le son como la gota de agua que cae del cubo, y como menudo polvo en las balanzas le son estimadas; he aquí que hace desaparecer las islas como polvo. Ni el Líbano bastará para el fuego, ni todos sus animales para el sacrificio. Como nada son todas las naciones delante de él; y en su comparación serán estimadas en menos que nada, y que lo que no es. ¿A qué, pues, haréis semejante a Dios, o qué imagen le compondréis?
\end{verso}

El profeta no rebaja a las naciones para despreciarlas, sino para medir a Dios: todo lo que existe pesa ante él lo que una gota que cae del cubo, y aun así es él quien lo sostiene. Pink condensa el cuadro en una fórmula: Dios ``lo sostiene todo, pero es independiente de todo. Él da a todos, pero de ninguno se enriquece'' (Pink, 2020, cap. 1). Pablo le da forma de doxología al final de su primera carta a Timoteo: el ``solo Soberano'', ``el único que tiene inmortalidad'' (1~Ti 6:15--16).

\precision{Sal 16:2 se ha leído como palabras de Cristo según la RV1909 (``mi bien a ti no aprovecha''; así Pink, 2020, cap. 1). La RVR1960 traduce el hemistiquio de otro modo (``No hay para mí bien fuera de ti''), y el hebreo \heb{טוֹבָתִי בַּל־עָלֶיךָ} es difícil. Por eso no se usa aquí como prueba.}

\subsection{Un conocimiento anterior a todo}

Si Dios es antes de todo, también lo son su conocimiento y su consejo: ``anuncio lo por venir desde el principio\ldots{} Mi consejo permanecerá'' (Is 46:10). Cuando el Nuevo Testamento habla de ese conocimiento previo, su objeto no son acciones sino \emph{personas}, como observa Pink (2020, cap. 4). Cristo fue ``entregado por el determinado consejo y anticipado conocimiento de Dios'' (Hch 2:23); Dios predestinó ``a los que antes conoció'' (Ro 8:29); no desechó ``a su pueblo, al cual desde antes conoció'' (Ro 11:2); y los creyentes son ``elegidos según la presciencia de Dios Padre'' (1~P 1:2). Es lo mismo que Dios dijo a Jeremías: ``Antes que te formase en el vientre te conocí'' (Jer 1:5). El Dios que existe antes de todo conoce a las personas antes de que existan, y ese conocimiento no es mera previsión, sino que descansa en su propósito: ``Dios previamente conoce lo que será porque Él ha decretado lo que será'' (Pink, 2020, cap. 4).

\precision{la relación entre presciencia y elección es un debate abierto entre tradiciones teológicas; la de Pink es la lectura reformada. Lo que este capítulo toma es la observación textual: el Dios que existe antes de todo conoce \emph{personas} antes de que existan. Al Dios omnisciente se le dedicará su propio capítulo.}

%===================================================================
\section{Un encuentro en la historia de Israel: los dioses que hay que cargar}

En el siglo~VI a.\,e.\,c., Judá está en el exilio de Babilonia. Los dioses del imperio parecen victoriosos: Bel y Nebo tienen templos, sacerdotes y procesiones. Isaías describe el día en que esos dioses tienen que huir:

\begin{verso}{Isaías 46:1–2}
Se postró Bel, se abatió Nebo; sus imágenes fueron puestas sobre bestias, sobre animales de carga; esas cosas que vosotros solíais llevar son alzadas cual carga, sobre las bestias cansadas. Fueron humillados, fueron abatidos juntamente; no pudieron escaparse de la carga, sino que tuvieron ellos mismos que ir en cautiverio.
\end{verso}

Los dioses de Babilonia dependen de todo: del platero que los hace (46:6), de los hombros que los cargan, de las bestias que se cansan. Si nadie los lleva, no se mueven (46:7). En contraste, el Dios de Israel dice:

\hebverso{וְעַד־זִקְנָה אֲנִי הוּא וְעַד־שֵׂיבָה אֲנִי אֶסְבֹּל \\ אֲנִי עָשִׂיתִי וַאֲנִי אֶשָּׂא וַאֲנִי אֶסְבֹּל וַאֲמַלֵּט}
\begin{verso}{Isaías 46:3–4}
Oídme, oh casa de Jacob, y todo el resto de la casa de Israel, los que sois traídos por mí desde el vientre, los que sois llevados desde la matriz. Y hasta la vejez yo mismo, y hasta las canas os soportaré yo; yo hice, yo llevaré, yo soportaré y guardaré.
\end{verso}

\begin{filologia}{Is 46:1--4}
\heb{נְשֻׂאֹתֵיכֶם} \emph{nesu'otejem} & \heb{נשׂא} & ``lo que vosotros cargabais'' (46:1): los ídolos son carga. \\
\heb{הַנְּשֻׂאִים מִנִּי־רָחַם} & \heb{נשׂא} & ``los llevados desde la matriz'' (46:3): Israel es cargado por Dios. El mismo verbo, en dirección contraria. \\
\heb{אֲנִי הוּא} \emph{ani hu} & --- & ``yo mismo'' (46:4): el Dios que no cambia con los años de su pueblo. \\
\heb{אֶסְבֹּל} \emph{esbol} & \heb{סבל} & ``soportaré'': llevar un peso a cuestas, como un cargador. \\
\end{filologia}

El texto opone dos religiones. En una, el ser humano carga a su dios. En la otra, Dios carga al ser humano desde que nace hasta que envejece. El Dios que no depende de nadie es el único que puede sostener a todos. Luego Isaías une la autosuficiencia con la presciencia:

\hebverso{מַגִּיד מֵרֵאשִׁית אַחֲרִית וּמִקֶּדֶם אֲשֶׁר לֹא־נַעֲשׂוּ \\ אֹמֵר עֲצָתִי תָקוּם וְכָל־חֶפְצִי אֶעֱשֶׂה}
\begin{verso}{Isaías 46:9–10}
Acordaos de las cosas pasadas desde los tiempos antiguos; porque yo soy Dios, y no hay otro Dios, y nada hay semejante a mí, que anuncio lo por venir desde el principio, y desde la antigüedad lo que aún no era hecho; que digo: Mi consejo permanecerá, y haré todo lo que quiero.
\end{verso}

``Desde el principio'' (\heb{מֵרֵאשִׁית}, \emph{mereshit}) es la palabra de Gn 1:1 (\emph{bereshit}). Dios anunció el final desde el principio porque estaba antes del principio. En el versículo siguiente llama ``desde el oriente al ave'' (46:11): un rey extranjero que libraría a los cautivos. La historia confirmó la palabra: Babilonia cayó y los exiliados volvieron.

\precision{bel es el título de Marduk, dios principal de Babilonia, y Nebo (Nabu) era su hijo; sus imágenes salían en procesión en las fiestas del año nuevo. El ``ave'' del oriente se identifica con Ciro de Persia (Is 44:28; 45:1), que tomó Babilonia en 539 a.\,e.\,c. (Oswalt, 1998). La lectura de Is 46 como ``encuentro'' que revela la preexistencia es propia de esta serie.}

%===================================================================
\section{La preexistencia en el judaísmo}\label{sec:pre-judaismo}

\subsection{``Solo el Santo y su Nombre''}
El midrash \emph{Pirqé de Rabí Eliézer} comienza su relato de la creación con una frase que resume este capítulo:

\hebverso{עַד שֶׁלֹּא נִבְרָא הָעוֹלָם הָיָה הַקָּדוֹשׁ בָּרוּךְ הוּא וּשְׁמוֹ הַגָּדוֹל בִּלְבַד}

``Antes de que el mundo fuese creado, existían solo el Santo, bendito sea, y su gran Nombre'' (\emph{Pirqé de Rabí Eliézer} 3). Y agrega que ``subió a su pensamiento crear el mundo'': la creación nace de una decisión, no de una necesidad.

\precision{interpretación rabínica (midrash de época medieval temprana), según Sefaria. Que el Nombre estuviera ``con'' Dios antes del mundo no significa, en el midrash, un segundo ser: es la forma de decir que Dios ya era quien se revelaría a Moisés.}

\subsection{``Yo fui, yo soy, yo seré''}
El \emph{Midrash Shemot Rabá} (3:6) recoge varias lecturas de Éx 3:14. Rabí Yitsjac explica por qué \emph{ehyeh} aparece tres veces en el versículo: Dios dijo a Moisés que dijera a Israel

\hebverso{אֲנִי שֶׁהָיִיתִי וַאֲנִי הוּא עַכְשָׁו וַאֲנִי הוּא לֶעָתִיד לָבֹא}

``yo soy el que fui, y yo soy ahora, y yo soy en el porvenir''. El mismo pasaje, en nombre de Rabí Abba bar Mamal, dice que Dios es llamado según sus obras (\heb{לְפִי מַעֲשַׂי אֲנִי נִקְרָא}): juez cuando juzga, \emph{Elohim}; misericordioso cuando se compadece, \SENOR. El Nombre de Éx 3:14 reúne todos los tiempos y todas las obras.

\subsection{``Estaré con ellos''}
El Talmud lee Éx 3:14 como promesa de presencia. Dios dijo a Moisés: ``Ve y diles a Israel: yo estuve con vosotros en esta servidumbre, y estaré con vosotros en la servidumbre de los reinos''. Moisés responde que basta la angustia de cada hora, y Dios le dice que les diga solo ``\emph{Ehyeh} me envió'' (\emph{Berajot} 9b). Rashi resume: \heb{אֶהְיֶה עִמָּם בְּצָרָה זוֹ}, ``estaré con ellos en esta angustia'' (Rashi a Éx 3:14).

\precision{interpretaciones rabínicas, según Sefaria (Talmud de Babilonia, edición William Davidson; Rashi). Para la tradición judía, el Nombre de la zarza es menos una definición del ser de Dios que una promesa de compañía. Esta lectura coincide con Éx 3:12 (``yo estaré contigo'') y equilibra la lectura de la Septuaginta.}

\subsection{El Targum: ``el que es y el que fue''}
Los targumim palestinos (traducciones arameas leídas en la sinagoga) amplían el Nombre con los tiempos del verbo. El Targum Pseudo-Jonatán traduce Éx 3:14 como ``el que es y el que será'', y en Dt 32:39 hace decir a Dios: ``Yo soy el que es y el que fue, y yo soy el que será''. Ese mismo lenguaje reaparece en el Apocalipsis: ``el que es y que era y que ha de venir'' (Ap 1:4, 8).

\precision{paráfrasis targúmicas citadas según Bauckham (1993), que propone que la fórmula del Apocalipsis depende de esta interpretación judía del Nombre y no de la filosofía griega. La fecha de redacción del Pseudo-Jonatán es tardía, aunque recoge tradiciones antiguas; conviene cotejar el texto arameo en Sefaria.}

\subsection{Maimónides: el primer existente}
La \emph{Mishné Torá} de Maimónides (s.~XII) comienza con este atributo:

\hebverso{יְסוֹד הַיְסוֹדוֹת וְעַמּוּד הַחָכְמוֹת לֵידַע שֶׁיֵּשׁ שָׁם מָצוּי רִאשׁוֹן}

``El fundamento de los fundamentos y el pilar de las ciencias es saber que hay un Primer Existente'' (\emph{Yesodé haTorá} 1:1). Si se supone que él no existe, nada podría existir (1:2). Y si se supone que nada existe fuera de él, él seguiría existiendo,

\hebverso{שֶׁכָּל הַנִּמְצָאִים צְרִיכִין לוֹ וְהוּא בָּרוּךְ הוּא אֵינוֹ צָרִיךְ לָהֶם}

``porque todos los seres lo necesitan a él, y él, bendito sea, no los necesita a ellos'' (1:3). Es la definición de este capítulo dicha por el mayor sabio medieval del judaísmo. En 1:4 concluye que por eso su verdad no es como la de ningún otro (cf. Jer 10:10, en el capítulo sobre la verdad).

\subsection{Lo que existía antes del mundo}
El Talmud enseña que siete cosas fueron creadas antes que el mundo: la Torá, el arrepentimiento, el jardín del Edén, la Guehena, el trono de gloria, el templo y \heb{שְׁמוֹ שֶׁל מָשִׁיחַ}, ``el nombre del Mesías'', según el salmo: ``Será su nombre para siempre, se perpetuará su nombre mientras dure el sol'' (\heb{לִפְנֵי־שֶׁמֶשׁ יִנּוֹן שְׁמוֹ}, Sal 72:17; \emph{Pesajim} 54a).

\precision{interpretación rabínica. La frase hebrea \emph{lifné shémesh} puede leerse ``mientras dure el sol'' (RVR1960) o ``antes del sol''; los sabios la leen ``antes''. Para el Talmud lo que precede al mundo es el \emph{nombre} del Mesías, es decir, el propósito de Dios de enviarlo, no una persona preexistente. La fe cristiana afirma más: el Mesías mismo existía con el Padre (Jn 17:5). El paralelo muestra que la idea de un plan de Dios anterior a la creación que incluye al Mesías es judía antes de ser cristiana.}

\subsection{La oración: \emph{Adon Olam}}
El himno \emph{Adon Olam}, que muchas comunidades cantan al comenzar o al cerrar la oración de la mañana, empieza así:

\hebverso{אֲדוֹן עוֹלָם אֲשֶׁר מָלַךְ בְּטֶרֶם כָּל־יְצִיר נִבְרָא}

``Señor del mundo, que reinó antes de que toda criatura fuera creada''. Más adelante dice \heb{וְהוּא הָיָה וְהוּא הֹוֶה וְהוּא יִהְיֶה}, ``y él fue, y él es, y él será'', y lo llama \heb{בְּלִי רֵאשִׁית בְּלִי תַכְלִית}, ``sin principio, sin fin''. La palabra \emph{béterem} de Sal 90:2 y los tres tiempos de Shemot Rabá entraron en la oración diaria. El himno termina confiando el alma a ese Dios al dormir: el que reinaba antes de todo cuida al que duerme.

\precision{el autor del \emph{Adon Olam} es desconocido; se lo atribuye a veces a Salomón ibn Gabirol, sin pruebas firmes. Entró en la liturgia de la mañana hacia el siglo~XV, y hay una costumbre de recitarlo antes de dormir y en el lecho de muerte. Texto según el sidur (Sefaria).}

\subsection{Humildad ante el ``antes''}
La Misná advierte sobre la curiosidad sin límite: ``Todo el que especula sobre cuatro cosas, mejor le habría sido no venir al mundo: qué hay arriba, qué hay abajo, qué hubo antes y qué habrá después'' (\emph{Jaguigá} 2:1). El midrash observa que la Torá empieza con la letra \emph{bet} (\heb{בְּרֵאשִׁית}), cerrada por arriba, por abajo y por detrás, y abierta solo hacia adelante: no se puede indagar lo que hubo antes, sino desde el día en que el mundo fue creado (\emph{Bereshit Rabá} 1:10).

\begin{verso}{Deuteronomio 29:29}
Las cosas secretas pertenecen a \SENOR{} nuestro Dios; mas las reveladas son para nosotros y para nuestros hijos para siempre, para que cumplamos todas las palabras de esta ley.
\end{verso}

\precision{interpretaciones rabínicas (Misná y midrash), según Sefaria; la cita de \emph{Bereshit Rabá} 1:10 conviene cotejarla en el original. Coinciden con el espíritu de esta obra: Dios solo puede conocerse en lo que él revela (1~Co 2:10--11).}

\subsection{Qumrán: todo estaba en su designio}
La \emph{Regla de la Comunidad} de Qumrán enseña que todo procede del ``Dios del conocimiento'':

\hebverso{מאל הדעות כול הויה ונהייה ולפני היותם הכין כול מחשבתם}

``Del Dios del conocimiento procede todo lo que es y lo que será; y antes de que existieran, estableció todo su designio'' (1QS 3:15). Es la preexistencia unida a la presciencia, con el mismo verbo \emph{hayá} de Éx 3:14: Dios es antes de lo que ``es'' (\heb{הויה}) y de lo que ``será'' (\heb{נהייה}).

\precision{traducción según García Martínez (1992) y Vermes (2004); conviene cotejar la imagen de 1QS en la Leon Levy Dead Sea Scrolls Digital Library. Los textos de Qumrán se citan sin vocales, como están en el manuscrito.}

%===================================================================
\section{La preexistencia en Cristo, leída con Pablo}

\subsection{``Antes que Abraham fuese, yo soy''}\label{sec:pre-abraham}
En el templo, durante una discusión sobre quiénes son los verdaderos hijos de Abraham, Jesús dice que Abraham ``se gozó de que había de ver mi día''. Sus oyentes objetan que no tiene cincuenta años. Jesús responde:

\begin{verso}{Juan 8:58–59}
Jesús les dijo: De cierto, de cierto os digo: Antes que Abraham fuese, yo soy. Tomaron entonces piedras para arrojárselas; pero Jesús se escondió y salió del templo; y atravesando por en medio de ellos, se fue.
\end{verso}

\begin{filologia}{Jn 8:58}
\grc{πρὶν Ἀβραὰμ γενέσθαι} & \grc{γίνομαι} & ``antes que Abraham llegara a ser'': aoristo de un verbo que significa comenzar a existir. Abraham tuvo un comienzo. \\
\grc{ἐγὼ εἰμί} & \grc{εἰμί} & ``yo soy'': presente del verbo ``ser'', sin predicado. Jesús no dice ``yo era'', sino ``yo soy'', como el \emph{ani hu} de Is 43:10 en la LXX. \\
\grc{ἦραν οὖν λίθους} & \grc{αἴρω} & ``tomaron entonces piedras'': la pena de la blasfemia contra el Nombre (Lv 24:16). \\
\end{filologia}

El contraste de verbos es el centro de la frase: Abraham \emph{llegó a ser}; Jesús \emph{es}. Sus oyentes entienden que habla en el lenguaje del Nombre, y por eso toman piedras. En el mismo capítulo Jesús había dicho ``si no creéis que yo soy, en vuestros pecados moriréis'' (8:24), y ``cuando hayáis levantado al Hijo del Hombre, entonces conoceréis que yo soy'' (8:28). En Getsemaní, al decir ``Yo soy'', los soldados ``retrocedieron, y cayeron a tierra'' (18:6).

\precision{la relación entre el \grc{ἐγὼ εἰμί} absoluto de Juan y el \emph{ani hu} de Isaías (LXX) es reconocida en los comentarios (Carson, 1991). Juan no es Pablo; se cita porque es el texto que propone el material de la serie y porque Pablo afirma lo mismo con otro lenguaje.}

\subsection{La comunión antes de la creación}
Juan comienza su evangelio con las palabras de Gn 1:1 y dice qué había ``en el principio'':

\begin{verso}{Juan 1:1–3}
En el principio era el Verbo, y el Verbo era con Dios, y el Verbo era Dios. Este era en el principio con Dios. Todas las cosas por él fueron hechas, y sin él nada de lo que ha sido hecho, fue hecho.
\end{verso}

\begin{filologia}{Jn 1:1--3, 18}
\grc{ἐν ἀρχῇ ἦν} & \grc{εἰμί} & ``en el principio era'': imperfecto, existencia continua. Cuando empezó el principio, el Verbo ya era. \\
\grc{πρὸς τὸν θεόν} & \grc{πρός} & ``con Dios'': la preposición sugiere orientación, ``vuelto hacia Dios'', cara a cara. \\
\grc{ἐγένετο} & \grc{γίνομαι} & ``fueron hechas'' (1:3): de las cosas se dice \emph{llegar a ser}; del Verbo, \emph{ser}. El mismo contraste de Jn 8:58. \\
\grc{εἰς τὸν κόλπον τοῦ πατρός} & \grc{κόλπος} & ``en el seno del Padre'' (1:18): la intimidad del hijo recostado en el pecho del padre. \\
\end{filologia}

Antes de que existiera el mundo, no había un Dios aislado: había comunión. Jesús lo dice en su oración al Padre:

\begin{verso}{Juan 17:5}
Ahora pues, Padre, glorifícame tú al lado tuyo, con aquella gloria que tuve contigo antes que el mundo fuese.
\end{verso}
\begin{verso}{Juan 17:24}
Padre, aquellos que me has dado, quiero que donde yo estoy, también ellos estén conmigo, para que vean mi gloria que me has dado; porque me has amado desde antes de la fundación del mundo.
\end{verso}

\grc{ᾗ εἶχον πρὸ τοῦ τὸν κόσμον εἶναι παρὰ σοί}: ``la gloria que tenía antes que el mundo existiera, junto a ti''. \grc{ἠγάπησάς με πρὸ καταβολῆς κόσμου}: ``me amaste antes de la fundación del mundo''. Antes de la creación había gloria compartida y amor entre el Padre y el Hijo. Por eso la autosuficiencia de Dios no es la de un ser solitario que crea para tener compañía: Dios no creó porque le faltara amor, sino porque el amor que ya existía entre el Padre y el Hijo quiso tener hijos (Jn 17:23--24; Ef 1:4--5).

\precision{esta conclusión es teología cristiana clásica: Dios es amor (1~Jn 4:8) desde antes de crear porque es Padre, Hijo y Espíritu, y por eso no necesitaba al mundo para amar (Grudem, 2007). El judaísmo no comparte esta lectura trinitaria; afirma con la misma fuerza que Dios no necesita al mundo (Maimónides, \emph{Yesodé haTorá} 1:3).}

\subsection{La Sabiduría junto a Dios}
Proverbios 8 presenta a la Sabiduría hablando de sí misma antes de la creación:

\hebverso{וָאֶהְיֶה אֶצְלוֹ אָמוֹן וָאֶהְיֶה שַׁעֲשֻׁעִים יוֹם יוֹם מְשַׂחֶקֶת לְפָנָיו בְּכָל־עֵת}
\begin{verso}{Proverbios 8:27, 30}
Cuando formaba los cielos, allí estaba yo; Cuando trazaba el círculo sobre la faz del abismo\ldots{} Con él estaba yo ordenándolo todo, Y era su delicia de día en día, Teniendo solaz delante de él en todo tiempo.
\end{verso}

La Sabiduría dice \heb{וָאֶהְיֶה אֶצְלוֹ} (\emph{va'ehyeh etsló}), ``y yo estaba junto a él'', con la forma \emph{ehyeh} de Éx 3:14, y describe ese ``antes'' como deleite y gozo delante de Dios. Pablo llama a Cristo ``sabiduría de Dios'' (1~Co 1:24, 30), y Col 1:15--17 usa el lenguaje de la Sabiduría de Proverbios 8.

\precision{proverbios 8 personifica la sabiduría como un atributo de Dios; no es en su contexto una profecía directa del Mesías. La lectura cristológica es tipológica y fue tratada en el capítulo sobre la sabiduría. Los sabios judíos la leen como la Torá, que Dios ``miraba'' al crear el mundo (\emph{Bereshit Rabá} 1:1). La expresión ``me poseía'' (\heb{קָנָנִי}, 8:22) se discute (``me adquirió'', ``me creó'', ``me engendró'') y fue centro de debates cristológicos antiguos.}

\subsection{Pablo: Cristo antes de todas las cosas}\label{sec:pre-pablo}
Pablo no usa la palabra ``preexistencia'', pero la supone en sus himnos y en sus fórmulas sobre el Hijo:

\begin{verso}{Colosenses 1:15–17}
El es la imagen del Dios invisible, el primogénito de toda creación. Porque en él fueron creadas todas las cosas, las que hay en los cielos y las que hay en la tierra, visibles e invisibles; sean tronos, sean dominios, sean principados, sean potestades; todo fue creado por medio de él y para él. Y él es antes de todas las cosas, y todas las cosas en él subsisten.
\end{verso}

\begin{filologia}{Col 1:15--17}
\grc{πρωτότοκος πάσης κτίσεως} & \grc{πρωτότοκος} & ``primogénito de toda creación'': rango y herencia del primogénito (Sal 89:27 [89:28]), no primera criatura; el v. 16 explica ``porque en él fueron creadas''. \\
\grc{αὐτός ἐστιν πρὸ πάντων} & \grc{πρό} & ``él es antes de todas las cosas'': presente, como el ``yo soy'' de Jn 8:58. \\
\grc{συνέστηκεν} & \grc{συνίστημι} & ``subsisten, se mantienen unidas'': perfecto. Todo depende de él para seguir existiendo; él no depende de nada. \\
\end{filologia}

Lo que Isaías decía del \SENOR{} (``el primero'', 44:6) y lo que Maimónides diría del Primer Existente (todo lo necesita, él no necesita nada), Pablo lo dice de Cristo: es antes de todo y todo subsiste en él. Otros textos suponen la misma realidad:

\begin{itemize}
\item \textbf{``Siendo en forma de Dios''} (\grc{ἐν μορφῇ θεοῦ ὑπάρχων}, Flp 2:6): el Hijo existía en la condición de Dios antes de tomar ``forma de siervo'' (2:7).
\item \textbf{``Siendo rico''} (\grc{πλούσιος ὤν}, 2~Co 8:9): ``por amor a vosotros se hizo pobre''. Solo puede empobrecerse quien antes era rico.
\item \textbf{``Dios envió a su Hijo''} (Gá 4:4; Ro 8:3): el Hijo es enviado ``cuando vino el cumplimiento del tiempo''; existía antes de nacer ``de mujer''.
\item \textbf{``La roca era Cristo''} (1~Co 10:4): Pablo lo ve acompañando a Israel en el desierto.
\end{itemize}

\begin{verso}{1 Corintios 8:6}
para nosotros, sin embargo, sólo hay un Dios, el Padre, del cual proceden todas las cosas, y nosotros somos para él; y un Señor, Jesucristo, por medio del cual son todas las cosas, y nosotros por medio de él.
\end{verso}

En 1~Co 8:6 Pablo toma el \emph{Shemá} (Dt 6:4: ``el \SENOR{} nuestro Dios, el \SENOR{} uno es'') y lo reparte entre ``un Dios, el Padre'' y ``un Señor, Jesucristo''. Del Padre todas las cosas proceden; por medio del Señor todas las cosas existen. Pablo pone a Jesús del lado del Creador, antes de toda criatura.

\precision{la lectura de 1~Co 8:6 como un \emph{Shemá} que incluye a Jesús en la identidad del único Dios es de Bauckham (2008) y Wright (2013). Dunn (1998) es más cauto respecto de la preexistencia en algunos textos (por ejemplo, Gá 4:4, donde ``enviar'' podría decirse también de un profeta), pero la reconoce en la tradición que Pablo recibe. La autoría de Colosenses se discute; Flp 2:6, 2~Co 8:9 y 1~Co 8:6 están en cartas indiscutidas.}

\subsection{Pablo en el Areópago: el Dios que no necesita nada}
Pablo predica a los atenienses, que tienen templos y altares para todos los dioses, y les anuncia al ``Dios no conocido'':

\begin{verso}{Hechos 17:24–25, 28}
El Dios que hizo el mundo y todas las cosas que en él hay, siendo Señor del cielo y de la tierra, no habita en templos hechos por manos humanas, ni es honrado por manos de hombres, como si necesitase de algo; pues él es quien da a todos vida y aliento y todas las cosas\ldots{} Porque en él vivimos, y nos movemos, y somos; como algunos de vuestros propios poetas también han dicho: Porque linaje suyo somos.
\end{verso}

\begin{filologia}{Hch 17:25}
\grc{θεραπεύεται} & \grc{θεραπεύω} & ``es servido, atendido'': la palabra del cuidado de los dioses en el templo (vestirlos, alimentarlos). \\
\grc{προσδεόμενός τινος} & \grc{προσδέομαι} & ``como si necesitase algo \emph{además}'': única vez en el NT. A Dios no le falta nada que se le pueda añadir. \\
\grc{αὐτὸς διδοὺς πᾶσι ζωήν} & \grc{δίδωμι} & ``él es quien da a todos vida'': la dirección es siempre de Dios hacia la criatura. \\
\end{filologia}

Es el sermón de Isaías 46 predicado a los griegos: los ídolos necesitan manos humanas; el Dios verdadero da las manos, la vida y el aliento. Pablo, judío formado en las Escrituras, habla aquí como el profeta del exilio.

\subsection{``¿Quién le dio a él primero?''}
Pablo cierra los capítulos 9--11 de Romanos, su gran reflexión sobre Israel, con una doxología hecha de citas del Antiguo Testamento:

\begin{verso}{Romanos 11:33–36}
¡Oh profundidad de las riquezas de la sabiduría y de la ciencia de Dios! ¡Cuán insondables son sus juicios, e inescrutables sus caminos! Porque ¿quién entendió la mente del Señor? ¿O quién fue su consejero? ¿O quién le dio a él primero, para que le fuese recompensado? Porque de él, y por él, y para él, son todas las cosas. A él sea la gloria por los siglos. Amén.
\end{verso}

\begin{center}\small\renewcommand{\arraystretch}{1.2}
\begin{tabular}{@{}L{2.6cm}L{11.8cm}@{}}
\toprule
Job 41:11 [41:3] TM & \heb{מִי הִקְדִּימַנִי וַאֲשַׁלֵּם}: ``¿quién se me adelantó, para que yo le pague?'' \\
Job 41:3 LXX & \grc{ἢ τίς ἀντιστήσεταί μοι καὶ ὑπομενεῖ;}: ``¿o quién se me opondrá y resistirá?'' (otro sentido) \\
Ro 11:35 & \grc{ἢ τίς προέδωκεν αὐτῷ, καὶ ἀνταποδοθήσεται αὐτῷ;}: ``¿o quién le dio primero, y le será devuelto?'' \\
Ro 11:34 & ``¿quién entendió la mente del Señor? ¿O quién fue su consejero?'' = Is 40:13 LXX \\
\bottomrule
\end{tabular}\par
{\footnotesize\color{gris}\textit{Tabla 4.} Job 41 en el texto masorético, la Septuaginta y Romanos.}
\end{center}

En Ro 11:35 Pablo sigue el sentido del hebreo y no el de la Septuaginta: \grc{προέδωκεν} (``dio antes'') traduce bien \heb{הִקְדִּימַנִי} (``se me adelantó''). Nadie llega antes que Dios. La conclusión \grc{ἐξ αὐτοῦ καὶ δι᾽ αὐτοῦ καὶ εἰς αὐτὸν τὰ πάντα} recorre todo: todas las cosas vienen \emph{de} él (origen), existen \emph{por} él (sostén) y van \emph{hacia} él (fin). Es la aseidad en una línea, convertida en alabanza.

\precision{que Ro 11:35 depende del texto hebreo de Job y Ro 11:34 de Is 40:13 LXX es la posición de los comentarios (Moo, 1996). La lectura de la doxología como síntesis de la aseidad es de esta serie, en la línea de Pink (2020, cap. 1).}

\subsection{Elegidos en él antes de la fundación del mundo}
La presciencia de Dios, en Pablo, tiene un centro: Cristo.

\begin{verso}{Efesios 1:3–5}
Bendito sea el Dios y Padre de nuestro Señor Jesucristo, que nos bendijo con toda bendición espiritual en los lugares celestiales en Cristo, según nos escogió en él antes de la fundación del mundo, para que fuésemos santos y sin mancha delante de él, en amor habiéndonos predestinado para ser adoptados hijos suyos por medio de Jesucristo, según el puro afecto de su voluntad.
\end{verso}
\begin{verso}{2 Timoteo 1:9}
quien nos salvó y llamó con llamamiento santo, no conforme a nuestras obras, sino según el propósito suyo y la gracia que nos fue dada en Cristo Jesús antes de los tiempos de los siglos.
\end{verso}

``Antes de la fundación del mundo'' (\grc{πρὸ καταβολῆς κόσμου}) aparece tres veces en el Nuevo Testamento, y las tres hablan de Cristo: el Padre lo amó (Jn 17:24), nos escogió en él (Ef 1:4) y él fue ``destinado'' (\grc{προεγνωσμένου}, literalmente ``conocido de antemano'') como Cordero (1~P 1:20). El amor del Padre al Hijo antes de la creación es el lugar donde Dios pensó en nosotros.

\precision{la unión de Jn 17:24, Ef 1:4 y 1~P 1:20 por la fórmula común es una lectura propia de esta serie. La autoría de Efesios y de 2~Timoteo se discute; Ro 8:29 (``a los que antes conoció'') afirma la presciencia en una carta indiscutida. 1~Pedro no es de Pablo, y se cita por la fórmula.}

\subsection{Vida en sí mismo}
\begin{verso}{Juan 5:26}
Porque como el Padre tiene vida en sí mismo, así también ha dado al Hijo el tener vida en sí mismo.
\end{verso}

\grc{ζωὴν ἔχει ἐν ἑαυτῷ}: ``tiene vida en sí mismo''. Es la definición de la aseidad dicha por Jesús. Ninguna criatura tiene vida en sí misma: la recibe y la pierde (Ec 12:7). El Padre la tiene, y la ha dado al Hijo para que también la tenga ``en sí mismo''. Por eso Jesús puede decir ``Yo soy la resurrección y la vida'' (Jn 11:25), y Pablo puede llamar a Cristo ``nuestra vida'' (Col 3:4). Lo que Éx 3:14 reveló en la zarza, Jn 1:14 lo muestra en carne: ``vimos su gloria''.

%===================================================================
\section{Aplicación: un atributo incomunicable}

La preexistencia es solo de Dios. Ninguna criatura existe antes de todo, por sí misma ni sin necesitar nada. Lo que este atributo produce en el creyente no es imitación, sino una forma de vivir delante de él.

\begin{enumerate}
\item \textbf{Adorar sin creer que le hacemos un favor.} Dios no es ``honrado por manos de hombres, como si necesitase de algo'' (Hch 17:25). La adoración no le suma nada; nos devuelve a nuestro lugar. ``Siervos inútiles somos, pues lo que debíamos hacer, hicimos'' (Lc 17:10).
\item \textbf{Depender sin vergüenza.} ``En él vivimos, y nos movemos, y somos'' (Hch 17:28). La autosuficiencia es de Dios; la del ser humano es una ilusión. Jesús lo dijo a sus discípulos: ``separados de mí nada podéis hacer'' (Jn 15:5).
\item \textbf{Dejar los dioses que hay que cargar.} Todo lo que necesita ser sostenido por nosotros para seguir en pie (dinero, prestigio, control) funciona como Bel y Nebo (Is 46:1--2). El Dios verdadero dice: ``yo hice, yo llevaré, yo soportaré y guardaré'' (46:4).
\item \textbf{Descansar en lo que nos precede.} Antes de que existiéramos, Dios ya nos conocía (Jer 1:5) y nos escogió en Cristo (Ef 1:4). La vida del creyente no empezó en su decisión, sino en el amor del Padre al Hijo ``desde antes de la fundación del mundo'' (Jn 17:24).
\item \textbf{Conocer con humildad.} ``Las cosas secretas pertenecen a \SENOR{} nuestro Dios'' (Dt 29:29). Del ``antes'' de Dios sabemos lo que él reveló, y eso basta para confiar.
\end{enumerate}

La pregunta que deja el atributo es personal: ¿cargo a mis dioses, o dejo que Dios me cargue?

%===================================================================
\section{Síntesis}

\begin{tcolorbox}[colback=papiro,colframe=azul,boxrule=1pt,arc=3pt]
Antes de que existiera algo, Dios era: ``desde el siglo y hasta el siglo, tú eres Dios'' (Sal 90:2). Existe por sí mismo (``YO SOY EL QUE SOY'', Éx 3:14), no depende de nada (Sal 50:12) y nadie le dio primero (Job 41:11). Su conocimiento y su designio son anteriores a la creación (Is 46:10; Jer 1:5). Israel lo confesó como \emph{ani hu} frente a los dioses que hay que cargar; el judaísmo lo llamó el Primer Existente. Jesús dijo ``antes que Abraham fuese, yo soy'' (Jn 8:58), y Pablo lo vio ``antes de todas las cosas'' (Col 1:17): antes de la creación, el Padre y el Hijo ya vivían en gloria y amor (Jn 17:5, 24).
\end{tcolorbox}

\begin{center}\small\renewcommand{\arraystretch}{1.3}
\begin{tabular}{@{}L{3.0cm}L{11.4cm}@{}}
\toprule
\textbf{Mirada} & \textbf{Aporte al atributo} \\ \midrule
Hebrea & \emph{Ehyeh asher ehyeh} (Éx 3:14); \emph{ani hu}; \emph{béterem}, \emph{miqqédem}; ``el primero y el postrero''; Job 41:11; Is 46: Dios carga, los ídolos son cargados. \\
Judía & ``Solo el Santo y su Nombre'' (PRE 3); ``yo fui, yo soy, yo seré'' (Shemot Rabá 3:6); \emph{Ehyeh} como presencia (Berajot 9b; Rashi); el Primer Existente (Maimónides); \emph{Adon Olam}; 1QS 3:15. \\
Teológica & Antes de crear, Dios era completo y sin necesidad; creó libremente; nada le añade; su presciencia mira a personas y descansa en su propósito. \\
Paulina & Cristo antes de todas las cosas (Col 1:17); ``siendo rico'' (2~Co 8:9); 1~Co 8:6; Hch 17:25; Ro 11:35--36; elegidos en él antes de la fundación del mundo (Ef 1:4). \\
\bottomrule
\end{tabular}\end{center}

\begin{verso}{Jeremías 9:24 [9:23]}
Mas alábese en esto el que se hubiere de alabar: en entenderme y conocerme, que yo soy \SENOR{}, que hago misericordia, juicio y justicia en la tierra; porque estas cosas quiero, dice \SENOR{}.
\end{verso}

``Que yo soy \SENOR'': conocer a Dios empieza por saber que él es antes que nosotros y sin nosotros, y que aun así decidió darse a conocer.

%===================================================================
\section*{Compendio bíblico}
\begin{center}\small\renewcommand{\arraystretch}{1.12}
\begin{tabular}{@{}L{2.2cm}L{5.0cm}|L{2.2cm}L{5.0cm}@{}}
\toprule
Gn 1:1 & ``En el principio creó Dios'' & Jn 1:1--3 & ``En el principio era el Verbo'' \\
Éx 3:14 & ``YO SOY EL QUE SOY'' & Jn 1:18 & ``que está en el seno del Padre'' \\
Dt 32:39 & ``Ved ahora que yo, yo soy'' & Jn 5:26 & ``tiene vida en sí mismo'' \\
Dt 33:27 & ``El eterno Dios es tu refugio'' & Jn 8:58 & ``Antes que Abraham fuese, yo soy'' \\
Neh 9:5--6 & ``desde la eternidad hasta la eternidad'' & Jn 17:5 & ``la gloria que tuve contigo'' \\
Job 41:11 & ``¿Quién me ha dado a mí primero?'' & Jn 17:24 & ``me has amado desde antes'' \\
Sal 50:12 & ``Si yo tuviese hambre'' & Hch 17:24--25 & ``como si necesitase de algo'' \\
Sal 90:2 & ``tú eres Dios'' & Ro 11:35--36 & ``¿quién le dio a él primero?'' \\
Sal 93:2 & ``Tú eres eternamente'' & 1~Co 8:6 & ``un Señor, Jesucristo'' \\
Pr 8:22--31 & ``allí estaba yo'' & 2~Co 8:9 & ``siendo rico'' \\
Is 40:15--18 & ``como la gota de agua'' & Flp 2:6 & ``siendo en forma de Dios'' \\
Is 43:10 & ``antes de mí no fue formado dios'' & Col 1:17 & ``él es antes de todas las cosas'' \\
Is 44:6 & ``Yo soy el primero'' & Ef 1:4 & ``antes de la fundación del mundo'' \\
Is 46:4 & ``yo llevaré, yo soportaré'' & 2~Ti 1:9 & ``antes de los tiempos de los siglos'' \\
Is 46:10 & ``anuncio lo por venir desde el principio'' & 1~P 1:20 & ``ya destinado desde antes'' \\
Jer 1:5 & ``Antes que te formase\ldots{} te conocí'' & Ap 1:8 & ``el que es y que era'' \\
Mi 5:2 & ``sus salidas son desde el principio'' & Ap 4:11 & ``por tu voluntad existen'' \\
Hab 1:12 & ``¿No eres tú desde el principio?'' & Ap 22:13 & ``el primero y el último'' \\
\bottomrule
\end{tabular}\par
{\footnotesize\color{gris}\textit{Tabla 5.} Compendio bíblico de la preexistencia. \textit{Nota.} Textos de la RVR1960.}
\end{center}

%===================================================================
\section*{Glosario del capítulo}
\begin{description}\setlength{\itemsep}{2pt}
\item[\emph{Adon Olam}] (\heb{אֲדוֹן עוֹלָם}). ``Señor del mundo'': himno litúrgico que confiesa al Dios que reinó antes de toda criatura.
\item[\emph{Ani hu}] (\heb{אֲנִי הוּא}). ``Yo (soy) él'': fórmula de Dios en Dt 32:39 e Isaías 40--48; en la LXX, \grc{ἐγώ εἰμι}.
\item[Aseidad] (latín \emph{a se}). Existencia por sí mismo, sin causa ni dependencia.
\item[\emph{Béterem}] (\heb{בְּטֶרֶם}). ``Antes que'' (Sal 90:2; Jer 1:5).
\item[\emph{Ego eimí}] (\grc{ἐγώ εἰμι}). ``Yo soy''; absoluto en Jn 8:58.
\item[\emph{Ehyeh}] (\heb{אֶהְיֶה}). ``Soy, seré, estaré'': forma del verbo \emph{hayá} en Éx 3:12, 14.
\item[\emph{Hayá}] (\heb{הָיָה}). Ser, existir, llegar a ser, estar presente.
\item[\emph{Hiqdim}] (\heb{הִקְדִּים}). Adelantarse, dar primero (Job 41:11).
\item[\emph{Kólpos}] (\grc{κόλπος}). Seno, regazo: ``en el seno del Padre'' (Jn 1:18).
\item[\emph{Matsuí rishón}] (\heb{מָצוּי רִאשׁוֹן}). ``Primer Existente'' (Maimónides).
\item[\emph{Miqqédem}] (\heb{מִקֶּדֶם}). Desde el principio, desde la antigüedad (Hab 1:12).
\item[\emph{Pro katabolés kósmou}] (\grc{πρὸ καταβολῆς κόσμου}). ``Antes de la fundación del mundo'' (Jn 17:24; Ef 1:4; 1~P 1:20).
\item[\emph{Prognósis}] (\grc{πρόγνωσις}). Presciencia, ``anticipado conocimiento'' (Hch 2:23; 1~P 1:2).
\item[\emph{Prosdéomai}] (\grc{προσδέομαι}). Necesitar algo además (Hch 17:25).
\item[\emph{Qédem}] (\heb{קֶדֶם}). Lo que está delante: el oriente; lo antiguo.
\item[\emph{Rishón}] (\heb{רִאשׁוֹן}). Primero (Is 44:6).
\item[\emph{Synéstēken}] (\grc{συνέστηκεν}). ``Subsisten, se mantienen unidas'' (Col 1:17).
\item[\emph{Yatsar}] (\heb{יָצַר}). Formar, como el alfarero (Gn 2:7; Is 43:10; Jer 1:5).
\end{description}

